\section{Conclusion}
\label{sec:conclusion}

Chiyoda reframes evacuation as an information-control problem. In a
hazard-coupled crowd, emergency communication changes beliefs, beliefs change
routes, and route convergence changes exposure and congestion. The framework
therefore measures not only whether information spreads, but whether spreading
it at a particular time, place, credibility, and radius improves safety.

The current empirical package supports that framing, but it also narrows the
claim. Static beacons are the strongest current baseline by
information-safety efficiency, global broadcast shows why reach is not enough,
and entropy-targeted intervention exposes how belief improvement can coexist
with weak aggregate completion. The next step is not to claim operational
readiness, but to test external validity across layouts, hazards, and
population mixes while preserving the safety-control interpretation developed
here.
